\subsection{纤维空间和群扩张的平凡性}

命题 2. --- 设 X 是局部紧空间，局部紧群 H 通过 \((x,\xi)\mapsto x\xi\) 在其中连续且 逆紧地从右作用。假设 X/H 仿紧。设 g 是 H 到 \(R^{n}\) 的连续表示。则存在从 X 到 \(R^{n}\) 的连续映射 f，使得 \(f(x\xi)=f(x)+g(\xi)\) 对所有 \(x\in X\) 和 \(\xi\in H\) 成立。

立即可归结为 n = 1 的情形。由于加法群 R 与乘法群 \(R_{+}^{*}\) 同构，此时命题立即由 §2, No.~4 的 Prop. 7 得出。

推论。--- 设 X 是局部紧空间，一个有限维实向量空间 V 通过 \((x,v)\mapsto xv\) 在其中连续且 逆紧地从右作用。令 \(\pi\) 为从 X 到 B = X/V 的典范映射。假设 B 仿紧。

\begin{enumerate}
\def\labelenumi{\alph{enumi})}
\item
  存在连续映射 \(f\)，定义在 \(\mathbf{X}\) 上取值于 \(\mathbf{V}\)，满足 \(f(xv) = f(x) + v\)，对所有 \(x \in \mathbf{X}\) 和 \(v \in \mathbf{V}\) 成立。
\item
  若 \(f\) 是满足 a) 条件的映射，则映射 \(x \mapsto (\pi(x), f(x))\) 是从 \(X\) 到 \(B \times V\) 的同胚。
\end{enumerate}

断言 a) 由 Prop. 2 得出，其中 g 取为 V 的恒等映射。设 f 是满足 a) 条件的映射。从 X 到 X 的映射 \(x \mapsto x \cdot (-f(x))\) 连续，并且在每个轨道上为常值，因而具有 \(\varphi \circ \pi\) 的形式，其中 \(\varphi\) 是从 B 到 X 的连续映射；对于每个 \(b \in B\)，有 \(\pi(\varphi(b)) = b\)。从 X 到 B × V 的映射 \(x \mapsto (\pi(x), f(x))\) 与从 B × V 到 X 的映射 \((b, v) \mapsto \varphi(b) \cdot v\) 互为逆映射，因为 \(\varphi(\pi(x)) \cdot f(x) = x \cdot (-f(x)) \cdot (f(x)) = x\)、\(\pi(\varphi(b) \cdot v) = \pi(\varphi(b)) = b\)，并且若 \(b = \pi(y)\)，则

\[
f (\varphi (\pi (y)) \cdot v) = f (y \cdot (- f (y)) \cdot v) = f (y) - f (y) + v = v.
\]

由于这些映射连续，它们是同胚。

注记。--- 设 E 是有限维实仿射空间，T 是紧空间，\(\mu\) 是 T 上总质量为 1 的测度，f 是从 T 到 E 的连续映射。若在 E 中选定原点 a，则 E 由此具有向量空间结构，因而积分 \(\int_{\mathrm{T}} f(t) d\mu(t)\) 有意义；它表示 E 中满足下式的点 x：

\[
x - a = \int_ {\mathrm{T}} (f (t) - a) d \mu (t).
\]

这个点与 \(a\) 的选择无关。事实上，设 \(a' \in \mathbf{E}\) 和 \(x' \in \mathbf{E}\) 满足 \(x' - a' = \int_{\mathbf{T}} (f(t) - a') d\mu(t)\)。则

\[
\begin{array}{c} x ^ {\prime} - a = (x ^ {\prime} - a ^ {\prime}) + (a ^ {\prime} - a) = \int_ {\mathrm{T}} \bigl (f (t) - a ^ {\prime} \bigr) d \mu (t) + \int_ {\mathrm{T}} (a ^ {\prime} - a) d \mu (t) \\ = \int_ {\mathrm{T}} \bigl (f (t) - a \bigr) d \mu (t) = x - a, \end{array}
\]

从而 \(x' = x\)。因此，可以使用符号 \(\int_{\mathrm{T}} f(t) d\mu(t)\)，而不必指明 E 中原点的选择。若 u 是从 E 到另一个有限维仿射空间 \(E'\) 的仿射映射，则

\[
u \Big (\int_ {\mathrm{T}} f (t) d \mu (t) \Big) = \int_ {\mathrm{T}} u \big (f (t) \big) d \mu (t).
\]

事实上，可以将 E 和 \(E'\) 以某种方式识别为向量空间，使 u 成为线性映射；此时该公式已知（Ch. III, §3, No.~2, Prop. 2 和 No.~3, Prop. 7）。

引理 2. --- 设 G 是紧群，\(\mu\) 是 G 的归一化 Haar 测度，E 是有限维实仿射空间，A 是 E 的仿射群，\(\rho\) 是从 G 到 A 的同态。假设对于每个 \(x \in \mathbf{E}\)，从 G 到 E 的映射 \(s \mapsto \rho(s)x\) 连续。则对于每个 \(x \in \mathbf{E}\)，点

\[
x _ {0} = \int_ {\mathrm{G}} \rho (s) x d \mu (s) \in \mathrm{E}
\]

在 G 下不变。

事实上，对于每个 \(t \in G\)，

\[
\rho (t) x _ {0} = \int_ {\mathrm{G}} \rho (t) \rho (s) x d \mu (s) = \int_ {\mathrm{G}} \rho (t s) x d \mu (s) = \int_ {\mathrm{G}} \rho (s) x d \mu (s) = x _ {0}.
\]

命题 3. --- 设 G 是局部紧群。设 H 是 G 的闭正规子群，同构于 \(\mathbf{R}^n\)，且 \(\mathbf{G} / \mathbf{H}\) 是紧的。

\begin{enumerate}
\def\labelenumi{\alph{enumi})}
\item
  存在闭子群 \(\mathbf{L}\)，它是 \(\mathbf{G}\) 的闭子群，使得 \(\mathbf{G}\) 是 \(\mathbf{L}\) 与 \(\mathbf{H}\) 的拓扑半直积。
\item
  若 \(\mathbf{M}\) 是 \(\mathbf{G}\) 的紧子群，则存在 \(x \in \mathbf{H}\)，使得 \(x^{-1}\mathbf{M}x \subset \mathbf{L}\)。
\item
  \(\mathbf{G}\) 的每个紧子群都包含于某个极大紧子群中。
\item
  \(\mathbf{G}\) 的极大紧子群正是 \(\mathbf{L}\) 在 \(\mathbf{G}\) 的内自同构下的变换像。
\end{enumerate}

令 \(\pi\) 为从 \(\mathbf{G}\) 到 \(\mathbf{K} = \mathbf{G} / \mathbf{H}\) 的典范同态。通过取商，映射 \((s,h)\mapsto shs^{-1}\) 从 \(\mathbf{G}\times \mathbf{H}\) 到 \(\mathbf{H}\) 定义了连续映射 \((\sigma ,h)\mapsto \sigma \cdot h\) 从 \(\mathbf{K}\times \mathbf{H}\) 到 \(\mathbf{H}\)，使得 \(shs^{-1} = \pi (s)\cdot h\)。我们将 \(\mathbf{H}\) 识别为 \(\mathbf{R}^n\)（因此，对于群 \(\mathbf{H}\) 的运算，视情形使用乘法或加法记号）。由 Prop. 2 的推论，存在连续映射 \(f\)，从 \(\mathbf{G}\) 到 \(\mathbf{H}\)，使得 \(f(xh) = f(x) + h\)，对 \(x\in \mathbf{G}, h\in \mathbf{H}\) 成立。对于每个 \(x\in \mathbf{G}\)，令 \(p(x) = x\cdot (-f(x))\)，它只依赖于 \(x\) 关于 \(\mathbf{H}\) 的陪集。置

\[
\begin{array}{r l} \mathrm{F} (x, y) & = p (x y) ^ {- 1} p (x) p (y) = f (x y) y ^ {- 1} x ^ {- 1} x (- f (x)) y (- f (y)) \\ & = f (x y) [ y ^ {- 1} (- f (x)) y ] (- f (y)) \\ & = f (x y) - \pi (y) ^ {- 1} f (x) - f (y). \end{array}\tag{1}
\]

我们看到，若对于 G 中所有 x,y 都有 \(F(x,y) = 0\)，则 \(p(G) = L\) 是 G 的子群，并且与 H 的每个陪集恰好交于一点。由于 p 连续，G 于是是 L 与 H 的拓扑半直积（GT, III, §2, No.~10）。

现在，对于任意 \(h, h' \in H\)，

\[
\begin{array}{r l} F (x h, y h ^ {\prime}) & = f (x h y h ^ {\prime}) - \pi (y) ^ {- 1} f (x h) - f (y h ^ {\prime}) \\ & = f (x h y) + h ^ {\prime} - \pi (y) ^ {- 1} f (x) - \pi (y) ^ {- 1} h - f (y) - h ^ {\prime} \\ & = f \big (x y (\pi (y) ^ {- 1} h) \big) - \pi (y) ^ {- 1} f (x) - f (y) - \pi (y) ^ {- 1} h \\ & = f (x y) - \pi (y) ^ {- 1} f (x) - f (y) = F (x, y). \end{array}
\]

因此，F 通过取商定义了连续映射 \(\varphi\)，从 \(K \times K\) 到 H。

另一方面，对于任意 \(x, y, z\)，在 \(\mathbf{G}\) 中，有

\[
\begin{array}{r l} & {\mathrm{F} (z, x y) + \mathrm{F} (x, y) = f (z x y) - \pi (x y) ^ {- 1} f (z) - f (x y) + f (x y)} \\ & {\qquad - \pi (y) ^ {- 1} f (x) - f (y)} \\ & {\qquad = \pi (y) ^ {- 1} f (z x) - \pi (x y) ^ {- 1} f (z) - \pi (y) ^ {- 1} f (x) + f (z x y)} \\ & {\qquad - \pi (y) ^ {- 1} f (z x) - f (y)} \\ & {\qquad = \pi (y) ^ {- 1} \mathrm{F} (z, x) + \mathrm{F} (z x, y),} \end{array}
\]

因此，对于任意 \(x', y', z'\)，在 \(\mathbf{K}\) 中，

\[
- \varphi (x ^ {\prime}, y ^ {\prime}) = \varphi (z ^ {\prime}, x ^ {\prime} y ^ {\prime}) - y ^ {- 1} \varphi (z ^ {\prime}, x ^ {\prime}) - \varphi (z ^ {\prime} x ^ {\prime}, y ^ {\prime}).
\]

对 \(z'\) 使用 K 的归一化 Haar 测度 \(\alpha\) 积分。置 \(\psi(x') = \int \varphi(z', x') d\alpha(z')\)，则 \(\psi\) 是 K 上的连续函数，并且（注意到由 GT, VII, §2, No.~1, Prop. 1，K 在 \(R^{n}\) 上的运算保持 \(R^{n}\) 的向量空间结构）得到

\[
- \varphi (x ^ {\prime}, y ^ {\prime}) = \psi (x ^ {\prime} y ^ {\prime}) - y ^ {- 1} \psi (x ^ {\prime}) - \psi (y ^ {\prime}).
\]

换言之，置 \(k - \psi \circ \pi\)，这是 G 上的连续函数，

\[
- \mathrm{F} (x, y) = k (x y) - \pi (y) ^ {- 1} k (x) - k (y).\tag{2}
\]

比较 (1) 和 (2) 可见，若将 \(f\) 替换为连续函数 \(f + k\)（它仍满足性质 \(f(xh) = f(x) + h\)），则 \(F\) 被替换为 0；如前所见，这就完成了 a) 的证明。

对于每个 \(g \in \mathbf{G}\)，令 \(l_g\)（分别为 \(h_g\)）为 \(\mathbf{L}\)（分别为 \(\mathbf{H}\)）中满足 \(g = h_g l_g\) 的唯一元素。若 \(h_1 \in \mathbf{H}\) 且 \(g \in \mathbf{G}\)，则

\[
g h _ {1} = h _ {g} l _ {g} h _ {1} = h _ {g} (l _ {g} h _ {1} l _ {g} ^ {- 1}) l _ {g},
\]

于是 \(h_{gh_1} = h_g + l_g h_1 l_g^{-1}\)。对于每个 \(g \in \mathbf{G}\)，令 \(\psi_g\) 为由下式定义的从 \(\mathbf{H}\) 到自身的映射：

\[
\psi_ {g} (h _ {1}) = h _ {g} + l _ {g} h _ {1} l _ {g} ^ {- 1}.
\]

容易看出，从 G×H 到 H 的映射 \((g,h_{1})\mapsto\psi_{g}(h_{1})\) 连续，并使 H 成为 G 的一个齐性空间，其中原点的稳定子为 L。此外，当将 H 识别为 \(R^{n}\) 时，\(\psi_{g}\) 是 H 到自身的仿射映射。现在设 M 是 G 的紧子群；由引理 2，存在 \(x\in H\)，使得 \(\psi_{m}(x)=x\) 对所有 \(m\in M\) 成立。对于 \(y\in H\)，\(\psi_{y}\) 是向量为 y 的平移；因此对于每个 \(m\in M\)，\(\psi_{x^{-1}}\circ\psi_{m}\circ\psi_{x}\) 将 H 的原点变换到自身，从而 \(x^{-1}mx\in L\)。这证明了 \(x^{-1}Mx\subset L\)，因而得到 b)。

设 \(L'\) 是 \(G\) 的包含 \(L\) 的闭子群。则 \(L'\) 是 \(L\) 与 \(L' \cap H\) 的拓扑半直积。若 \(L'\) 是紧的，则 \(L' \cap H\) 紧，因而退化为一点（GT, IV, §2, No.~2, Cor. 1 of Th. 2），所以 \(L' = L\)。这证明了 \(L\) 是 \(G\) 的极大紧子群；因此，\(L\) 在 \(G\) 的内自同构下的变换像也具有同样性质。于是 Prop. 3 的 c) 和 d) 是 b) 的直接推论。

命题 4. --- 设 G 是局部紧群，H 是 G 的闭正规子群，且 K = G/H 是紧的。则 H 到 R 的任意连续表示 u，只要满足 \(u(s\xi s^{-1}) = u(\xi)\)，对所有 \(\xi \in H\) 和 \(s \in G\) 成立，都可以延拓为 G 到 R 的连续表示。

令 \(L = G \times R\)，令 M 为 \((\xi, -u(\xi))\) 所成的集合，其中 \(\xi\) 遍历 H。显然 M 是 L 的闭正规子群。令 \(L' = L/M\)，令 \(\pi\) 为从 L 到 \(L'\) 的典范映射。由 M 与 R 在 L 中生成的子群是 \(H \times R\)，因而是闭的；所以 \(\pi(\mathbf{R})\) 是 \(L'\) 的闭子群 N。限制 \(\rho\)（即 \(\pi\) 在 R 上的限制）是从 R 到 N 的双射连续表示。附录 1 的引理 2 证明 \(\rho\) 是双连续的。此外，\(L'/N\) 同构于 \(\mathrm{L}/(\mathrm{H} \times \mathbf{R}) = \mathrm{G}/\mathrm{H}\)，因而是紧的。由 Prop. 3，并注意到 N 位于 \(L'\) 的中心，\(L'\) 是 N 与另一个子群的乘积。因此存在从 \(L'\) 到 N 的连续表示，在 N 上化为恒等映射。于是存在从 L 到 R 的连续表示 v，它在 M 上为平凡表示，并在 R 上化为恒等映射。对于 \(\xi \in H\)，有 \(v((\xi,0)) = v((\xi,-u(\xi))(e,u(\xi))) = u(\xi)\)，命题得证。

引理 3. --- 设 G 是由 e 的某个紧邻域生成的拓扑群。设 H 是 G 的闭子群，且齐性空间 G/H 是紧的。则 H 由 H 中 e 的某个紧邻域生成。

取紧集 C，使得 G = CH。必要时扩大 C，可设 C 生成 G，且 G = \(\overset{\circ}{CH}\)。于是 \(C^{2}\) 是紧的，并被 \(\overset{\circ}{Cs}\)（\(s \in H\)）覆盖，这些集合是开集。因此存在 H 中的 \(s_{1}, \ldots, s_{n}\)，使得 \(C^{2} \subset \overset{\circ}{Cs}_{1} \cup \cdots \cup \overset{\circ}{Cs}_{n}\)。令 \(\Gamma\) 为 H 中由 \(s_{i}\) 生成的子群。则 \(C^{2} \subset \mathrm{C}\Gamma\)。归纳可得，对每个 n 都有 \(C^{n} \subset \mathrm{C}\Gamma\)，因而 G = C \(\Gamma\)。H 的每个元素都可写成 ab，其中 \(a \in C\)、\(b \in \Gamma\)，从而 \(a \in H\)，继而 \(a \in C \cap H\)。因此 H 由 \(C \cap H\) 和 \(s_{i}\) 生成，也就是由一个紧集生成。

引理 4. --- 设 G 是连通拓扑群，D 是 G 的全不连通正规子群。则 D 包含于 G 的中心中。

事实上，设 \(d \in \mathbf{D}\)。\(\mathbf{G}\) 在连续映射 \(x \mapsto x dx^{-1}\) 下的像是 \(\mathbf{D}\) 的连通子集，因而退化为 \(\{d\}\)，这就证明了 \(xd = dx\) 对所有 \(x \in \mathbf{G}\) 成立。

命题 5. --- 设 G 是具有离散正规子群 D 的连通拓扑群，且 K = G/D 是紧的，并且 K 的换位子群在 K 中稠密。则 D 是有限的，G 是紧的。

群 G 与 K 局部同构（GT, III, §2, No.~6, Prop. 19），因而是局部紧的；由于它连通，所以由 e 的某个紧邻域生成。由引理 3 和引理 4，D 是有限生成阿贝尔群，因而同构于群 \(Z^{r} \times D_{1}\)，其中 \(D_{1}\) 有限（A, VII, §4, No.~7, Th. 3）。假设 r \textgreater{} 0。则存在从 D 到 Z 的表示 f。由 Prop. 4，f 可延拓为 G 到 R 的连续表示 g。通过取商，g 定义了 K 到 R/Z 的连续表示 \(g'\)；由于 R/Z 是阿贝尔的，\(g'\) 的核包含 K 的换位子群，因而 \(g'\) 是平凡的；换言之，\(g(\mathbf{G}) \subset \mathbf{Z}\)。由于 G 连通，遂有 \(g(\mathbf{G}) = \{0\}\)，但这与 \(f(\mathbf{D}) = \mathbf{Z}\) 矛盾。因此 r = 0，D 是有限的。于是 G 是紧的（GT, III, §4, No.~1, Cor. 2 of Prop. 2）。
% semantic-fragments: legacy-input-seam
\relax{}
